\documentclass[tikz, border=5pt]{standalone}
\usepackage{amsmath}
\usepackage{amssymb}
\usepackage{ctex}

\usetikzlibrary{arrows.meta, calc}

\begin{document}
\begin{tikzpicture}[>=Latex, thick]

% --- 定义基础参数 ---
% 适当增大单位 1 的长度，让水平实轴更宽敞
\def\U{5} 
% 将辅助轴角度调小至 22 度，使其与 45 度的对角线完全错开
\def\auxAngle{22}

% --- 样式定义 ---
\tikzset{axis/.style={->, line width=0.8pt}}
\tikzset{construction/.style={solid, line width=0.5pt}}
\tikzset{tick/.style={line width=0.6pt}}
\tikzset{m_label/.style={font=\large}}
\tikzset{crisis/.style={red!75!black}} % 第一次数学危机相关的元素

% --- 绘制实数轴 (\mathbb{R}) ---
\draw[axis] (-0.5, 0) -- (\U*1.75, 0) node[right, m_label] {$\mathbb{R}$};
\node[below, m_label, yshift=-3pt] at (0,0) {$0$};

% 标记单位 1
\coordinate (P1_real) at (\U, 0); 
\node[below, m_label, yshift=-3pt] at (\U, -0.1) {$1$};
\fill (P1_real) circle (1.2pt);

% ==========================================
% 第一次数学危机：构造 \sqrt{2}
% ==========================================
% 1. 绘制边长为 1 的正方形 (虚线)
\draw[construction, dashed, crisis] (0,0) rectangle (\U, \U);

% 2. 绘制对角线
\draw[construction, thick, crisis] (0,0) -- (\U, \U) node[midway, above left, xshift=-2pt] {$\sqrt{2}$};

% 3. 绘制圆弧将对角线长度落到实轴
\draw[construction, thick, crisis, ->] (\U, \U) arc (45:0:{\U*sqrt(2)});

% 4. 标记实轴上的 \sqrt{2} 点
\coordinate (P_sqrt2) at ({\U*sqrt(2)}, 0);
\draw[tick, crisis] (P_sqrt2) ++(0,-0.1) -- ++(0,0.2);
\node[below, m_label, crisis, yshift=-3pt] at (P_sqrt2) {$\sqrt{2}$};
\fill[crisis] (P_sqrt2) circle (1.2pt);


% ==========================================
% 有理数构造部分
% ==========================================
% 辅助轴伸长，避免拥挤
\draw[axis] (0,0) -- (\auxAngle:8.2) node[right, m_label] {辅助轴};

% 定义辅助轴上的关键节点的绝对长度，拉开它们之间的间距
\def\Lone{1.5}
\def\Lm{4.5}
\def\Ln{6.5}

\coordinate (A_1) at (\auxAngle:\Lone);
\coordinate (A_m) at (\auxAngle:\Lm); 
\coordinate (A_n) at (\auxAngle:\Ln); 

% 构造平行线
% 1. n 到 1
\draw[construction] (A_n) -- (P1_real);
\node[above left, m_label, xshift=-1pt, yshift=2pt] at (A_n) {$n$};
\draw[tick] (A_n) ++({\auxAngle+90}:0.08) -- ++({\auxAngle-90}:0.16);

% 2. 1 到 1/n
\coordinate (P_1_over_n) at ($(0,0)!{\Lone/\Ln}!(P1_real)$); 
\draw[construction] (A_1) -- (P_1_over_n);
\draw[tick] (A_1) ++({\auxAngle+90}:0.08) -- ++({\auxAngle-90}:0.16);
\node[below, m_label, yshift=-3pt] at (P_1_over_n) {$\frac{1}{n}$};
\node[above left, m_label, xshift=-1pt, yshift=2pt] at (A_1) {$1$};

% 3. m 到 m/n
\coordinate (P_m_over_n) at ($(0,0)!{\Lm/\Ln}!(P1_real)$); 
\draw[construction] (A_m) -- (P_m_over_n);
\draw[tick] (A_m) ++({\auxAngle+90}:0.08) -- ++({\auxAngle-90}:0.16);
\node[below, m_label, yshift=-3pt] at (P_m_over_n) {$\frac{m}{n}$};
\node[above left, m_label, xshift=-1pt, yshift=2pt] at (A_m) {$m$};

% --- 强调关键交点 ---
\fill (0,0) circle (1.2pt);
\fill (A_n) circle (1.2pt);
\fill (A_1) circle (1.2pt);
\fill (A_m) circle (1.2pt);
\fill (P_1_over_n) circle (1.2pt);
\fill (P_m_over_n) circle (1.2pt);

\end{tikzpicture}
\end{document}